% !TEX program = pdflatex

%%%%%%%%%%%%%%%%%%%%
% affineConwayCompData.tex
%%%%%%%%%%%%%%%%%%%%%

\pdfoutput=1 %When you use pdflatex
\documentclass{amsart} %When you use pdflatex
%\documentclass[dvipdfmx]{amsart} %When you use uplatex + dvipdfmx


\usepackage{tikz}
\usetikzlibrary{calc}
\usetikzlibrary{decorations.markings}
\usetikzlibrary {arrows.meta}
\usetikzlibrary {bending}
\usetikzlibrary{intersections}

\usepackage{amssymb}
\usepackage{amsmath}
\usepackage{url}
\usepackage{enumerate}
\usepackage{graphicx}

\usepackage{hyperref}
\hypersetup{
 colorlinks = true,
 linkcolor = blue,
 filecolor = blue,
 urlcolor = blue,
 citecolor = blue,
}
\hypersetup{unicode=true}

%\newcommand{\secK3}{\texorpdfstring{$K3$}{K3}}
\renewcommand{\baselinestretch}{1.4}

\usepackage{bm}
%\usepackage{enumerate}
\usepackage{enumitem}

\input mymacros
\input mymacrosaffineConway


%%%%%%%%%%%%%%%%%%%%%%%%%%%%

\begin{document}

\title[[Involutions in the affine Conway group: Computational data]
{Involutions in the affine Conway group: Computational data}


\author{Ichiro Shimada}
\address{Department of Mathematics,
Graduate School of Science,
Hiroshima University,
1-3-1 Kagamiyama,
Higashi-Hiroshima,
739-8526 JAPAN}
\email{ichiro-shimada@hiroshima-u.ac.jp}
\thanks{Supported by JSPS KAKENHI Grant Number~20H00112, 23K20209, and 23H00081}

%超平面配置に関連する離散構造の拡張、深化とその応用
%23H00081
%代数幾何学の計算機による研究の新展開
%23K20209
%格子、保型形式とK3曲面、エンリケス曲面の研究
%20H00112
\maketitle
%
%
% 
In the  files
\begin{itemize}
\item[]  \tname{affineConwayCompdata.txt} (??KB), 
\end{itemize}
 we present  the computational data used in the preprint 
\begin{itemize}
\item []
Ichiro Shimada.  Involutions in the affine {C}onway group. Preprint.
  \end{itemize}
This document explains the structure and contents of the data.
%
\par
{\bf Remarks}
\begin{itemize}
\item
In the following, we freely use the notation introduced in the preprint above.
\item
We use the format of {\tt GAP}.
In particular, 
some parts of the data are given in  the {\tt Record} format of {\tt GAP}.
\item
Elements of a vector space are given by \emph{row} vectors.
A linear map is expressed by the multiplication of a matrix from the \emph{right}.
\end{itemize}
%
\section{Contents of \tname{affineConwayCompdata.txt}}
%
In the following, we arrange the elements of $\Omega=\PP^1(\FF_{23})$ in order as 
\[
\infty, 0, 1,2 , \dots, 22, 
\]
and we  fix isomorphisms  
$\FF_2^{\Omega}\cong \FF_2^{24}$ and  $\ZZ^{\Omega}\cong \ZZ^{24}$
accordingly.
For example, the vector $v_{\infty}\in \ZZ^{\Omega}$ is $[1, 0, \dots, 0]$
and $v_{22}\in \ZZ^{\Omega}$ is $[0, \dots, 0, 1]$.
%
\begin{itemize}
%
\item \tname{theOmega} is the list
$\tt{[infinity, 0,1,2,\dots, 22]}$.
%
%\item \tname{GolayVects} is the extended binary Golay code $\Golay\subset \FF_2^{24}$
%constructed by the method in~\cite[Chapter~10.2]{TheCSbook}.
%This data is the list of $2^{12}=4096$ codewords of $\Golay$.
%Each codeword is a vector of length $24$
%  with entries in $\{0, 1\}$.
  %
 \item \tname{GolayWords} is the extended binary Golay code $\Golay\subset \FF_2^{24}$
constructed by the method in~\cite[Chapter~10.2]{TheCSbook}.
This data is the list of $2^{12}=4096$ codewords of $\Golay$,
each of which 
is expressed as a subset of \tname{theOmega}.
%
%  \item \tname{GolayWords} is another expression of \tname{GolayVects}.
%Each codeword is expressed as a subset of \tname{theOmega}.
%
%\item \tname{BasisLeech} gives a basis of the Leech lattice $\Lambda\subset \ZZ^{\Omega}$
%constructed by the method in~\cite[Chapter~10.3]{TheCSbook}.
%This data \tname{BasisLeech}  is a list of $24$ vectors in $\ZZ^{\Omega}$.
\item \tname{LeechBasis} gives a basis of the Leech lattice $\Lambda\subset \ZZ^{\Omega}$
constructed by the method in~\cite[Chapter~10.3]{TheCSbook}.
This data  %\tname{LeechBasis}  
is a list of $24$ vectors in $\ZZ^{\Omega}$.
%
\item \tname{GramLeech} is the Gram matrix of  $\Lambda$
with respect to \tname{LeechBasis}.
%
\item \tname{erecs} is the list of $4$ records,
each of which describes the involution $\vep_n=\vep(C_n)\in \dotz$
for $n=8, 12, 16, 24$.
Each record has the following items.
%
%
\begin{itemize}
%
\item \tname{n} is $n\in \{8, 12, 16, 24\}$.
%
%
\item \tname{word} is the  element $C_n$ of \tname{GolayWords} such that $\vep_n=\vep(C_n)$.
%
\item \tname{invol} is the matrix representation of the involusion $\vep_n\in \dotz$
with respect to the basis  \tname{LeechBasis} of $\Leech$.
%
\item \tname{plusbasis} is a basis of $\Leech(\vep_n, +)$,
which is a list of $(24-n)$ vectors in $\Leech$,
each of which is written in terms of  the basis  \tname{LeechBasis}.
(When $n=24$, this item is the null list $[\phantom{a}]$.)
%
\item \tname{OGplusrec}  is a record describing $\OG(\Leech(\vep_n, +))$.
It consists of the following items:
\begin{itemize}
\item \tname{gens}   is a finite generating set
of $\OG(\Leech(\vep_n, +))$.
Each element is reprented by a matrix of size $24-n$
with respect   \tname{plusbasis}.
(When $n=24$, this item is the null list $[\phantom{a}]$.)
\item \tname{size}  is the size of  $\OG(\Leech(\vep_n, +))$.
(When $n=24$, the size is $1$.)
\end{itemize}
%
\item \tname{minusbasis} is a basis of $\Leech(\vep_n, -)$,
which is a list of $n$ vectors in $\Leech$. 
(When $n=24$, this item is the list of row vectors of $\Id_{24}$.)
%
\item \tname{OGminusrec}  is a record describing $\OG(\Leech(\vep_n, +))$.
It consists of the following items:
\begin{itemize}
\item \tname{gens}   is a finite generating set
of $\OG(\Leech(\vep_n, +))$.
Each element is reprented by a matrix of size $24-n$
with respect   \tname{minusbasis}.
\item \tname{size}  is the size of  $\OG(\Leech(\vep_n, +))$.
\end{itemize}
%
 \end{itemize}
%
%\item \tname{K16} is the kernel of the natural homomorphism $\OG(\Leech(\vep_{16}, -)) \to \OG(q_{16})$.
%This data \tname{K16}  is a  list of $512$  matrices
%expressing isometries of $\Leech(\vep_{16}, -)$ 
%given by the Gram matrix  \tname{erec16.minusOGrec.Gram},
%where \tname{erec16} is the third element of the list \tname{epsilonrecs}.
%
 \end{itemize}
%
%
\section{The calculation of $\OG(M)$ for a positive-definite lattice $M$}
%
Let $M$ be a positive-definite lattice.
The finite group $\OG(M)$ is described by the method of \emph{stabilizer chain}.
This method provides us with the order and a finite generating set of $\OG(M)$.
\par
Let $m$ be the rank of $M$.
We fix a basis $[e_1, \dots, e_m]$ of $M$.
Elements of $M$ are expressed by vectors 
with respect to this basis, and 
elements of  $\OG(M)$ are expressed by matrices
with respect to this basis.
We choose another  basis $[b_1, \dots, b_m]$ of $M$.
For the computation,
 it is better to choose $[b_1, \dots, b_m]$
 consisting of \emph{short}  vectors.
 For $i=0, \dots, m$, 
we put
\[
\Sigma_i:=\set{g\in \OG(M)}{b_1^g=b_1, \dots, b_i^g=b_i},
\]
and consider the chain 
\[
\OG(M)=\Sigma_0\supset \dots\supset \Sigma_m=\{\Id\}.
\]
Then we have a natural bijection between 
the orbit 
\[
b_i^{\Sigma_{i-1}}=\set{b_i^g}{g\in \Sigma_{i-1}}
\]
 of $b_i$ by 
the group $\Sigma_{i-1}$ and 
the set $\Sigma_i\backslash \Sigma_{i-1}$.
We choose a subset 
$R_i \subset \Sigma_{i-1}$ such that the  projection 
$\Sigma_{i-1} \to\Sigma_i\backslash \Sigma_{i-1}$ gives a bijection
\[
R_i\;\;\cong\;\;  \Sigma_i\backslash \Sigma_{i-1}.
\]
For $i=1$, since $-\Id\in \Sigma_0$,
we have a subset $R_1\sprime \subset R_1$ such that
\[
R_1=R_1\sprime\cup (-R_1\sprime),
\qquad
R_1\sprime\cap (-R_1\sprime)=\emptyset.
\]
Then each element 
$g\in \OG(M)$ 
is written uniquely as 
\[
g=r_m \cdots r_2 r_1\sprime \delta \qquad
 (r_m\in R_m,\;  \dots, \; r_2\in R_2, \;\; r_1\sprime\in R_1\sprime,\;\;  \delta\in \{\pm \Id\}).
\]
In particular,  the order of $\OG(M)$ is equal %to the product of sizes of these $R_i$,
\[
2\,|R_m|\cdot \cdots \cdot|R_2| \cdot |R_1\sprime|, 
\]
and $\OG(M)$ is generated by the union 
\[
R_m\, \cup\, \cdots\,  \cup \,R_2\,\cup\, R_1\sprime\,\cup\,  \{- \Id\}.
\]
\erase{
\par
The old \tname{OGPosLat} 
returns a recond whose items contains 
%
\begin{itemize}
\item \tname{Gram} is the Gram matrix of $M$ with respect to the basis $[e_1, \dots, e_m]$.
\item \tname{basis} is a  basis $[b_1, \dots, b_m]$ of $M$.
\item \tname{gens} is the list $[R_m, \dots, R_2,  R_1\sprime, \{\pm  \Id\}]$.
\item \tname{orbs}   is the list $[S\sprime_1, S_2, \dots, S_m]$,
where $S_1\sprime$ is the list  of the pairs $[b_1^g, g]$,
where $g$ runs through $R_1\sprime$, and  for $i=2, \dots, m$, 
the list $S_i$ consists of  pairs $[b_i^g, g]$, 
where $g$ runs through $R_i$.
\item \tname{order} is the order of $\OG(M)$.
\end{itemize}
%
}
\par
The new \tname{MakeOGrec} makes from  the result of \tname{OGPosLat}  
a record \tname{OMrec}.
The record \tname{OMrec} describing $\OG(M)$
has the following items.
%
\begin{itemize}
\item \tname{Gram} is the Gram matrix of $M$ with respect to the basis $[e_1, \dots, e_m]$.
\item \tname{bs} is a  basis $[b_1, \dots, b_m]$ of $M$.
\item \tname{Rs} is the list $[R_1\sprime, R_2, \dots, R_m]$ 
of sets $R_1\sprime$ and $R_i $ for $i=2, \dots, m$.
\item \tname{order} is the order of $\OG(M)$.
\end{itemize}

%
\bibliographystyle{plain}
\bibliography{myrefsaffineConway}
%
\end{document}


